\chapter{Shares, Cap Tables, SAFEs, and Priced Venture Rounds}
\label{ch:cap-tables-safes}
\chapterline

In Parts~I and~II we examined businesses through their financial statements and cash-flow valuations: how assets balance liabilities plus equity, how accrual earnings differ from cash flow, and how a mature company's free cash flows imply an enterprise value. Now step to the Day~0 side of the ledger. You and a colleague have built a working prototype, hold almost no cash, and need to decide how to finance the business, how to structure corporate shares, and how to raise seed or venture capital without losing sight of either cash survival or ownership math.

Before touching share certificates, we first map the financial trajectory of a young company -- why some businesses fund growth purely from customer receipts while others intentionally run years of accounting losses funded by outside equity. We then build the exact algebra of corporate shares, option pools, convertible SAFEs, and priced venture rounds on the \textbf{capitalization table} (or \textbf{cap table}).

\begin{historybox}
For centuries, long-distance maritime ventures were temporary partnerships. Merchants pooled coin to outfit a ship, split whatever cargo returned to port, and liquidated the partnership at the end of the voyage. If one partner needed cash mid-voyage, there was no clean exit short of liquidating physical assets.

In 1602, the Dutch Republic chartered the \emph{Vereenigde Oostindische Compagnie} (VOC, the Dutch East India Company) in Amsterdam with a structural rule that created the modern joint-stock corporation: capital paid into the company became \emph{permanent}. Investors could no longer demand their coin back from the company's treasury after a single voyage; instead, to convert their stake into cash, they sold transferable shares to another buyer on the Amsterdam bourse. Locking capital inside the corporate entity allowed multi-decade investments in fleets and infrastructure while giving shareholders secondary-market liquidity.

Three and a half centuries later, permanent corporate equity met high-failure-rate engineering. Commercial banks cannot lend money at $8\%$ interest to a startup with no collateral and an $80\%$ failure rate, because capped $8\%$ interest on the winners cannot cover principal lost on the failures. In 1946, Harvard Business School professor Georges Doriot founded the American Research and Development Corporation (ARDC), raising institutional funds to buy minority equity stakes in risky postwar technology companies -- including a \$70{,}000 equity investment in Digital Equipment Corporation (DEC) in 1957 that grew past \$350 million. That same year in California, Arthur Rock arranged equity financing for eight young physicists and engineers leaving Shockley Semiconductor Laboratory to found Fairchild Semiconductor, establishing the modern venture structure in which founders, early employees, and risk investors share a common equity ledger.

As cloud infrastructure lowered the cost of starting a software company in the 2000s, spending \$40{,}000 in legal fees to negotiate a full preferred-stock charter for a \$250{,}000 seed check became disproportionate. Seed investors turned first to convertible promissory notes, and in late 2013 Y Combinator partner Carolynn Levy introduced the \textbf{SAFE} (\emph{Simple Agreement for Future Equity}), removing debt maturity dates and interest accrual altogether. In 2018, Y Combinator released the \textbf{Post-Money SAFE}, redesigning the conversion denominator so that founders and seed investors could compute exact pre-Series-A ownership dilution the moment a seed check was signed.
\end{historybox}

\section{The Startup Financial Trajectory: Capital, Revenue, and Profitability}
\label{sec:startup-financial-trajectory}

Engineers looking at startups from the outside often puzzle over an apparent contradiction: how can a software company report \$15{,}000{,}000 in annual GAAP net losses while venture funds compete to value it at \$200{,}000{,}000, whereas another profitable software firm earning \$2{,}000{,}000 in cash profit raises zero outside capital? Resolving that puzzle starts with separating two distinct financing engines.

\subsection{Two Financing Engines: Customer Cash vs.\ the Venture J-Curve}

At time $t$ (measured in months from launch), suppose a company generates monthly revenue $R(t) \ge 0$, earns gross margin $m \in (0, 1]$ after variable cost of goods sold (hosting, payment-network fees, support, or fulfillment), and spends fixed monthly operating expenses $E_{\text{op}}(t) > 0$ on engineering, product, sales, and administration. Its monthly operating cash surplus before working-capital timing and capital expenditures is $m\,R(t) - E_{\text{op}}(t)$. Every new company runs on one of two capital models:

\begin{enumerate}
  \item \textbf{The Customer-Funded (Bootstrapped) Engine}: the founders keep fixed payroll and overhead $E_{\text{op}}(t)$ low enough that gross profit $m\,R(t)$ exceeds $E_{\text{op}}(t)$ within months (often accelerated by upfront annual customer prepayments, which create positive operating cash flow before accrual revenue is even recognized). Growth is initially bounded by internal cash generation, yet the founders suffer $0\%$ external dilution, grant zero preferred liquidation hurdles, retain $100\%$ board control, and can either distribute annual cash dividends forever or accept a \$15{,}000{,}000 acquisition that makes every founder wealthy.
  \item \textbf{The Venture-Backed J-Curve Engine}: in markets characterized by heavy upfront R\&D (new silicon, foundation models, payment-rail licensing, biotech) or strong scale and network effects, waiting to hire engineers out of slow initial customer cash flow risks losing the market entirely to a faster-moving competitor. The startup therefore sells minority equity stakes for outside cash $C_0$, intentionally setting $E_{\text{op}}(t) > m\,R(t)$ for several years. Cumulative cash flow traces a \textbf{J-curve}: plunging downward through a multi-year investment trough (negative GAAP Net Income and negative Operating Cash Flow) before compounding gross profit $m\,R(t)$ overtakes fixed overhead $E_{\text{op}}(t)$ and climbs steeply upward. Accepting venture capital trades ownership dilution and preferred exit hurdles for time compression.
\end{enumerate}

\subsection{Stage-by-Stage Financial Anatomy: Pre-Seed to Profitability}

Negative GAAP Net Income in an early venture-backed company can reflect two completely opposite economic realities:
\begin{itemize}
  \item \textbf{Healthy Investment Burn (Positive Unit Economics)}: each customer cohort generates high gross margin ($m = 70\%$ to $85\%$ in software) and pays back its customer acquisition cost (CAC) within 12 to 18 months, so existing cohorts are solidly profitable; GAAP Net Income is negative solely because the company is spending upfront R\&D and sales payroll today to build products and acquire \emph{next} year's cohorts.
  \item \textbf{Broken Unit Economics (Negative Contribution Margin)}: each transaction loses money on variable cost ($m \le 0$) or requires customer acquisition subsidies larger than lifetime gross profit ($\text{LTV} < \text{CAC}$), so growing revenue actively widens losses.
\end{itemize}

Table~\ref{tab:startup-stages} maps how capital raised, revenue, unit economics, and GAAP profitability evolve from Day~0 to maturity.

\begin{table}[htbp]
\centering
\caption{Financial anatomy of a venture-backed technology company across stages.}
\label{tab:startup-stages}
\footnotesize
\setlength{\tabcolsep}{4.2pt}
\begin{tabular}{@{}lp{2.05cm}p{2.25cm}p{2.75cm}p{3.15cm}@{}}
\toprule
\textbf{Stage} & \textbf{Typical Capital \& Instrument} & \textbf{Revenue Profile} & \textbf{Unit Economics \& GAAP Profit} & \textbf{Required Proof Point for Next Stage} \\
\midrule
Pre-Seed & \$100k--\$1M (Bootstrapped or SAFEs) & \$0 to early pilot fees & Unproven margin; small burn (lean founder team) & Working product + acute customer pain validated \\
\addlinespace[2pt]
Seed & \$1M--\$4M (Post-Money SAFEs or Priced Seed) & \$100k--\$1.5M ARR & Positive gross margin $m > 0$ on initial users; GAAP loss & Product-market fit: organic retention \& usage pull \\
\addlinespace[2pt]
Series A & \$5M--\$15M (Priced Preferred) & \$1.5M--\$5M ARR & Proven cohort payback; GAAP loss from R\&D + GTM build & Repeatable sales motion ($\text{LTV} > 3\times \text{CAC}$) \\
\addlinespace[2pt]
Series B+ & \$15M--\$60M+ (Priced Preferred) & \$8M--\$50M+ ARR & High $m$; improving operating-margin ratio as fixed OpEx scales & Efficient multi-product / geographic expansion \\
\addlinespace[2pt]
Maturity / Exit & Self-funding FCF, IPO, or M\&A & \$50M--\$500M+ revenue & Positive GAAP Net Income \& Free Cash Flow ($m R > E_{\text{op}}$) & Durable cash compounding independent of capital markets \\
\bottomrule
\end{tabular}
\end{table}

\subsection{Burn Rate, Cash Runway, and Default Alive vs.\ Default Dead}

Suppose a startup holds cash balance $C_0 > 0$ in its bank account, earns monthly revenue $R(t)$ at positive gross margin $m \in (0, 1]$, and incurs fixed monthly operating expenses $E_{\text{op}}(t)$. Its monthly \textbf{gross burn} is total monthly cash outflow ($\text{COGS}(t) + E_{\text{op}}(t) = (1 - m)R(t) + E_{\text{op}}(t)$), while its monthly \textbf{net burn} is the net cash drained from the bank account each month:
\begin{equation}\label{eq:net-burn}
  B(t) = E_{\text{op}}(t) - m\,R(t).
\end{equation}
When net burn is approximately constant at $B > 0$, static \textbf{cash runway} is $\tau_{\text{static}} = C_0 / B$ months.

Static runway understates survival when gross profit is growing, yet overstates safety if founders assume another investor will always wire a bailout check before $\tau_{\text{static}}$ hits zero. Essayist and Y Combinator co-founder Paul Graham posed the sharper operational test: assuming expenses stay disciplined at $E_0$ and gross profit grows organically at its current rate without raising another dollar, does the company reach profitability ($m\,R(t) \ge E_0$) \emph{before} running out of cash (\textbf{Default Alive}), or does cash hit zero first (\textbf{Default Dead})?

Let us derive the exact analytical threshold. Suppose fixed monthly operating expenses are held at $E_0 > 0$, initial monthly gross profit is $G_0 = m\,R_0 < E_0$ (so initial net burn is $B_0 = E_0 - G_0 > 0$), and gross profit compounds continuously at monthly growth rate $g > 0$:
\[
  G(t) = G_0\,e^{g t}.
\]
Monthly net burn $B(t) = E_0 - G_0 e^{g t}$ reaches zero at the \textbf{break-even month} $t^*$ satisfying $G_0 e^{g t^*} = E_0$:
\begin{equation}\label{eq:breakeven-time}
  t^* = \frac{1}{g}\,\ln\!\left(\frac{E_0}{G_0}\right) = \frac{1}{g}\,\ln\!\left(\frac{E_0}{m\,R_0}\right).
\end{equation}
Total cash consumed between $t = 0$ and break-even $t^*$ is the integral of net burn up to $t^*$ (using $G_0 e^{g t^*} = E_0$):
\begin{equation}\label{eq:default-alive-cash}
  \Delta C_{\text{req}} = \int_0^{t^*} \bigl(E_0 - G_0\,e^{g u}\bigr)\,\dd u
  = E_0\,t^* - \frac{G_0}{g}\bigl(e^{g t^*} - 1\bigr)
  = \frac{1}{g}\left[E_0\ln\!\left(\frac{E_0}{m\,R_0}\right) - (E_0 - m\,R_0)\right].
\end{equation}
The startup is \textbf{Default Alive} if and only if $C_0 > \Delta C_{\text{req}}$, leaving a positive minimum cash trough $C_{\min} = C_0 - \Delta C_{\text{req}} > 0$ at $t = t^*$, after which cash builds every month.

\begin{example}[Testing Default Alive on Live Numbers]
Suppose a seed startup holds $C_0 = \text{USD } 1{,}800{,}000$ in cash, spends fixed monthly payroll and overhead $E_0 = \text{USD } 200{,}000$, and generates $R_0 = \text{USD } 100{,}000$ in monthly SaaS revenue at $m = 80\%$ gross margin ($G_0 = \text{USD } 80{,}000$). Current net burn is $B_0 = 200{,}000 - 80{,}000 = \text{USD } 120{,}000$ per month, giving a static runway of $1{,}800{,}000 / 120{,}000 = 15.0$ months. Note that $\ln(E_0 / G_0) = \ln(2.5) \approx 0.9163$, so the numerator inside brackets in~\eqref{eq:default-alive-cash} is $200{,}000(0.9163) - 120{,}000 = \text{USD } 63{,}258$:
\begin{itemize}
  \item \textbf{At $g = 6\%$ monthly organic growth ($g = 0.06$)}: break-even arrives at $t^* = 0.9163 / 0.06 = 15.27$ months, and total cash burned to reach break-even is $\Delta C_{\text{req}} = 63{,}258 / 0.06 = 1{,}054{,}302$ dollars, well below \$1{,}800{,}000. Even though $t^* = 15.27$ exceeds the $15.0$-month static runway, shrinking monthly burn keeps cumulative cash consumed well below $C_0$. The company is \textbf{Default Alive} with \$745{,}698 in reserve at its cash trough -- able to negotiate Series~A from strength rather than desperation.
  \item \textbf{At $g = 2.5\%$ monthly organic growth ($g = 0.025$)}: break-even recedes to $t^* = 0.9163 / 0.025 = 36.65$ months, requiring $\Delta C_{\text{req}} = 63{,}258 / 0.025 = 2{,}530{,}325$ dollars, exceeding \$1{,}800{,}000. The company is \textbf{Default Dead} by \$730{,}325 unless management immediately trims $E_0$ or accelerates high-margin revenue growth.
\end{itemize}
\end{example}

\begin{trapbox}
\textbf{The Fundraising Treadmill and Negative Unit Margins.}
Look at equation~\eqref{eq:default-alive-cash} when gross margin $m \le 0$ (for example, delivering a \$12 meal or an unpriced AI inference workload for \$8 while paying \$6 in variable cost plus \$4 in customer promotional credits). When $m \le 0$, $m\,R(t) \le 0$ for every positive revenue level: there is no finite break-even time $t^*$ at \emph{any} growth rate $g$. Growing top-line revenue increases net burn $B(t) = E_{\text{op}}(t) + \abs{m}R(t)$ faster, trapping the company on a fundraising treadmill dependent on continuous external equity infusions. If venture markets tighten and management cuts marketing to conserve cash, top-line growth halts and negative contribution margins still prevent reaching break-even. Venture equity can rationally finance fixed R\&D and recoverable customer acquisition; it should never subsidize structurally negative unit economics without a concrete engineering mechanism that turns $m$ positive.
\end{trapbox}

\section{Corporate Share Anatomy}
\label{sec:corporate-share-anatomy}

\subsection{Why Corporate Law Tracks Integer Shares Rather Than Fixed Percentages}

Once founders incorporate and begin sharing ownership with co-founders, hires, and investors, why does corporate law track millions of integer shares rather than writing ownership percentages directly into contracts?

Suppose two engineers discuss splitting a startup in percentages: ``Alice gets $50\%$, Bob gets $50\%$, and when Carol joins next month as founding engineer we will give her $10\%$.'' Immediately $50\% + 50\% + 10\% = 110\%$. Worse, suppose an early advisor asks for a clause stating ``Advisor shall hold $5\%$ of the company on a non-dilutable basis.'' If one holder is contractually pinned at $5\%$ forever while the company raises three rounds of outside capital, $100\%$ of every future investor's dilution is forced onto the unprotected founders and employees, driving founder stakes down and causing institutional investors to refuse to invest until the non-dilutable promise is removed.

Corporate law eliminates percentage contradictions by treating percentages as \emph{derived views}, never stored primitives. The stored state in the corporate stock ledger is an integer count of discrete units called \textbf{shares}. If holder $i$ owns $s_i \in \Z_{\ge 0}$ shares out of a chosen total denominator $S = \sum_{j} s_j$, then holder $i$'s ownership fraction is
\[
  \omega_i = \frac{s_i}{S}.
\]
When the company issues $\Delta S$ new shares to Carol or to a venture investor, Alice's integer share count $s_{\text{Alice}}$ does not change by a single share. What changes is the denominator $S \mapsto S + \Delta S$, scaling every existing holder's fraction uniformly by $\frac{S}{S + \Delta S}$ so that $\sum_i \omega_i = 1$ holds by construction.

\begin{intuitionbox}
Think of a corporation's stock ledger as an append-only state machine in which each participant holds an integer counter $s_i$ and primary share issuances never mutate existing holders' counters. Issuing $\Delta S$ new shares appends one row to the ledger rather than rewriting every prior shareholder's contract, while preserving $\sum_i \omega_i = 1$ automatically.
\end{intuitionbox}

\subsection{The Four Share-Count Denominators}

Because ownership is always a ratio $s_i / S$, any statement that ``you own $4{,}000{,}000$ shares'' means nothing until you specify the denominator $S$. Every corporation tracks four distinct share totals, and confusing them produces six-figure miscalculations.

\begin{definition}[Capitalization Denominators]
At any fixed point in a corporation's life:
\begin{enumerate}
  \item \textbf{Authorized Shares} ($S_{\text{auth}}$): the statutory ceiling in the current Certificate of Incorporation on how many shares the board of directors may issue (for instance, $10{,}000{,}000$ shares of Common Stock at incorporation). Unissued authorized shares carry zero votes, receive zero economic payouts, and never enter an ownership denominator. Raising $S_{\text{auth}}$ requires a formal charter amendment approved by the board and stockholders.
  \item \textbf{Issued and Outstanding Shares} ($S_{\text{out}}$): shares actually sold or issued to named legal persons and not repurchased or cancelled (for example, $4{,}000{,}000$ shares issued to Alice and $4{,}000{,}000$ to Bob, so $S_{\text{out}} = 8{,}000{,}000$). Statutory stockholder votes and immediate legal ownership use $S_{\text{out}}$ (including preferred stock voting on an as-converted basis).
  \item \textbf{Reserved Employee Option Pool} ($S_{\text{pool}} = S_{\text{granted}} + S_{\text{unalloc}}$): authorized common shares set aside by the board under an Equity Incentive Plan for employees, advisors, and consultants -- split into \textbf{granted options} ($S_{\text{granted}}$, promised to specific team members under signed option grants) and \textbf{unallocated options} ($S_{\text{unalloc}}$, reserved for future hires not yet recruited).
  \item \textbf{Fully Diluted Shares} ($S_{\text{FD}}$): total shares assuming every outstanding convertible preferred share, warrant, granted option, and unallocated pool share has been converted or exercised into common stock:
  \[
    S_{\text{FD}} = S_{\text{common}} + S_{\text{pref, as-converted}} + S_{\text{granted}} + S_{\text{unalloc}} + S_{\text{warrants}}.
  \]
\end{enumerate}
\end{definition}

Suppose a startup authorizes $S_{\text{auth}} = 10{,}000{,}000$ common shares at incorporation, issues $4{,}000{,}000$ shares to each of two founders ($S_{\text{out}} = 8{,}000{,}000$), reserves $1{,}000{,}000$ shares in an employee option pool ($300{,}000$ granted to a first engineer and $700{,}000$ unallocated), and leaves $1{,}000{,}000$ authorized shares unissued and unreserved. Here $S_{\text{out}} = 8{,}000{,}000$ and $S_{\text{FD}} = 8{,}000{,}000 + 300{,}000 + 700{,}000 = 9{,}000{,}000$ (never $10{,}000{,}000$). Each founder holds $4{,}000{,}000 / 8{,}000{,}000 = 50.00\%$ of current voting stock, yet $4{,}000{,}000 / 9{,}000{,}000 = 44.44\%$ on a fully diluted basis. Venture investors price rounds on $S_{\text{FD}}$ because every reserved option is expected to be granted to hires who build the company.

\subsection{Par Value, 409A Fair Market Value, and Preferred Price}

A single startup routinely carries three distinct dollar prices on its shares on the same day:

\begin{itemize}
  \item \textbf{Par Value} (typically \$0.00001 per share in a Delaware technology corporation): a nominal statutory floor below which the corporation may not originally issue shares. Setting par value at \$0.00001 allows two founders buying $8{,}000{,}000$ initial common shares to wire $8{,}000{,}000 \times 0.00001 = 80.00$ dollars (\$80.00 in cash) into the corporate bank account so the shares are validly issued and fully paid without requiring thousands of dollars of personal cash on Day~1. (When filing annual Delaware franchise taxes with $10{,}000{,}000$ low-par shares, always elect the \emph{Assumed Par Value Capital Method}, which typically yields the \$400 minimum tax, rather than the default \emph{Authorized Shares Method}, which generates an alarming five-figure bill based solely on share count.)
  \item \textbf{Common Fair Market Value (409A Valuation)}: the appraised fair market value of a single share of \emph{Common Stock}, determined at least annually (and after each financing) by an independent valuation firm under Section~409A of the US Internal Revenue Code. This figure sets the minimum legal strike price for employee stock options so employees avoid immediate tax penalties.
  \item \textbf{Preferred Price Per Share}: the cash price paid by venture investors for newly issued \emph{Preferred Stock}. Because Preferred Stock carries downside liquidation preferences, anti-dilution adjustments, and governance vetoes that Common Stock lacks (as derived in Chapter~\ref{ch:founder-contracts}), an early-stage startup's 409A Common FMV is legitimately much lower -- typically $20\%$ to $35\%$ of the simultaneous Preferred Stock price.
\end{itemize}

\section{Pre-Money Valuation, Post-Money Valuation, and Dilution}
\label{sec:pre-post-money}

Suppose an existing startup has $S_{\text{pre}}$ fully diluted shares before raising capital. A venture fund agrees to wire cash $I > 0$ into the company's treasury (a \textbf{primary issuance}, meaning the corporation issues brand-new shares in exchange for cash deposited into the corporate bank account, rather than buying existing shares from the founders).

\begin{formalbox}
\textbf{Conservation Equations of a Priced Equity Round.}
Let $V_{\text{pre}} > 0$ denote the agreed \textbf{pre-money valuation} of the company immediately before the investment, $I > 0$ the cash wired into the company, and $S_{\text{pre}}$ the pre-round fully diluted share count. Immediately after closing, the company holds every asset and capability it held before plus $I$ in new cash, giving the \textbf{post-money valuation}
\begin{equation}\label{eq:post-money}
  V_{\text{post}} = V_{\text{pre}} + I.
\end{equation}
For the investor's post-round ownership fraction $\alpha$ to match the fraction of total post-money value contributed in cash, we require
\begin{equation}\label{eq:investor-fraction}
  \alpha = \frac{I}{V_{\text{post}}} = \frac{I}{V_{\text{pre}} + I}.
\end{equation}
Pricing existing shares at \textbf{price per share} $p$ and issuing $S_{\text{new}}$ new shares to the investor gives
\begin{equation}\label{eq:pps-and-new-shares}
  p = \frac{V_{\text{pre}}}{S_{\text{pre}}}, \qquad
  S_{\text{new}} = \frac{I}{p} = S_{\text{pre}}\,\frac{I}{V_{\text{pre}}}, \qquad
  S_{\text{post}} = S_{\text{pre}} + S_{\text{new}}.
\end{equation}
\end{formalbox}

Check immediately that equations~\eqref{eq:post-money}--\eqref{eq:pps-and-new-shares} satisfy both ownership conservation and value conservation:
\[
  \frac{S_{\text{new}}}{S_{\text{post}}}
  = \frac{S_{\text{pre}}(I / V_{\text{pre}})}{S_{\text{pre}}(1 + I / V_{\text{pre}})}
  = \frac{I}{V_{\text{pre}} + I}
  = \alpha,
  \qquad
  p\,S_{\text{post}} = p\,S_{\text{pre}} + p\,S_{\text{new}} = V_{\text{pre}} + I = V_{\text{post}}.
\]
Every pre-existing holder who owned $s_i$ shares before the round still owns $s_i$ shares afterward; their ownership fraction scales from $s_i / S_{\text{pre}}$ to $\frac{s_i}{S_{\text{pre}}}(1 - \alpha)$, while the implied paper value of their shares across the closing wire remains $p\,s_i = \frac{s_i}{S_{\text{pre}}}\,V_{\text{pre}} = \frac{s_i}{S_{\text{post}}}\,V_{\text{post}}$, compounding upward if the company deploys cash $I$ into high-return engineering and distribution.

\begin{example}[Why Shrinking Percentage Can Mean Rising Stake Value]
Suppose you and a co-founder start with $4{,}000{,}000$ shares each ($50.00\%$ each of an $8{,}000{,}000$-share seed concept valued at \$2{,}000{,}000, so your half is worth \$1{,}000{,}000). Across seed and Series~A rounds that put \$7{,}000{,}000 of engineering and sales capital onto the balance sheet, total fully diluted shares expand to $11{,}428{,}571$ and the company reaches a funded post-money valuation $V_{\text{post}}$ of \$25{,}000{,}000 ($p = \text{USD } 25{,}000{,}000 / 11{,}428{,}571 = \text{USD } 2.1875$ per share). Your $4{,}000{,}000$ shares now represent
\[
  \frac{4{,}000{,}000}{11{,}428{,}571} = 35.00\%
\]
of the company. Two distinct changes occurred at once:
\begin{enumerate}
  \item \textbf{Economic stake value rose}: at \$2.1875 per share, your $4{,}000{,}000$ shares carry an implied value of $0.35 \times 25{,}000{,}000 = 8{,}750{,}000$ dollars (\$8{,}750{,}000), up from \$1{,}000{,}000, because primary cash entered the company's treasury and funded value creation above the dilution suffered.
  \item \textbf{Governance voting share fell}: your voting weight dropped irreversibly from $50.00\%$ to $35.00\%$. Below $50\%$ you can no longer unilaterally elect a majority of a common-voted board or block ordinary stockholder resolutions on share count alone.
\end{enumerate}
Evaluating any financing requires tracking both numbers simultaneously: implied dollar value per founder share and post-round voting governance.
\end{example}

\section{The Pre-Money Option Pool Shuffle}
\label{sec:option-pool-shuffle}

Read almost any institutional venture term sheet and you will encounter a clause worded like this:

\begin{quote}
\emph{``Amount Raised: \$4{,}000{,}000 at a Pre-Money Valuation of \$16{,}000{,}000 (\$20{,}000{,}000 Post-Money), including an unallocated Post-Money Employee Option Pool equal to $10.0\%$ of the Post-Money Fully Diluted Capital.''}
\end{quote}

Read casually, that sentence sounds as though the investors are valuing the founders' existing shares at \$16{,}000{,}000. Read algebraically, they are pricing the founders' existing shares as if the pre-money valuation were only \textbf{\$14{,}000{,}000} -- a \$2{,}000{,}000 reduction known in venture practice as the \textbf{option pool shuffle}.

Look closely at the words \emph{Pre-Money} and \emph{Post-Money} inside that single sentence. The required option pool is sized as a fraction $\theta_{\text{pool}} = 10\%$ of \emph{post-money} shares, so its dollar value at closing is $\theta_{\text{pool}} V_{\text{post}} = 0.10 \times 20{,}000{,}000 = 2{,}000{,}000$ dollars (\$2{,}000{,}000). Yet the term sheet requires those new unallocated option shares to be inserted \emph{inside the pre-money share denominator} before computing the investor's entry price per share $p$. Let us prove the exact identity.

\begin{proposition}[Effective Valuation Under a Pre-Money Option Pool]
\label{prop:option-pool-shuffle}
Suppose existing shareholders hold $S_{\text{exist}}$ fully diluted shares (with zero existing unallocated options), and an investor invests cash $I$ at a stated pre-money valuation $V_{\text{pre}}$ (so $V_{\text{post}} = V_{\text{pre}} + I$), subject to creating a new unallocated option pool equal to fraction $\theta_{\text{pool}} \in (0, 1 - I/V_{\text{post}})$ of post-money fully diluted shares $S_{\text{post}}$ inside the pre-money capitalization. Then:
\begin{enumerate}
  \item The round price per share is
  \begin{equation}\label{eq:effective-pps}
    p = \frac{V_{\text{pre}} - \theta_{\text{pool}} V_{\text{post}}}{S_{\text{exist}}}.
  \end{equation}
  \item Existing shareholders' shares are priced at an \textbf{effective pre-money valuation} of
  \begin{equation}\label{eq:effective-pre-money}
    V_{\text{eff}} = p\,S_{\text{exist}} = V_{\text{pre}} - \theta_{\text{pool}} V_{\text{post}},
  \end{equation}
  pushing $100\%$ of the dilution from the newly created option pool onto existing shareholders prior to the investor's check.
\end{enumerate}
\end{proposition}

\begin{proof}
After closing, total fully diluted shares are $S_{\text{post}} = S_{\text{exist}} + S_{\text{pool}} + S_{\text{inv}}$ at uniform price per share $p$, satisfying $p\,S_{\text{post}} = V_{\text{post}}$. The investor's cash $I$ buys fraction $\alpha = I / V_{\text{post}}$ of $S_{\text{post}}$, and the new unallocated option pool occupies fraction $\theta_{\text{pool}}$ of $S_{\text{post}}$. Because the three blocks partition $S_{\text{post}}$, existing holders retain post-money fraction
\[
  \frac{S_{\text{exist}}}{S_{\text{post}}} = 1 - \alpha - \theta_{\text{pool}} = 1 - \frac{I}{V_{\text{post}}} - \theta_{\text{pool}} = \frac{V_{\text{post}} - I - \theta_{\text{pool}} V_{\text{post}}}{V_{\text{post}}} = \frac{V_{\text{pre}} - \theta_{\text{pool}} V_{\text{post}}}{V_{\text{post}}}.
\]
Multiplying both sides by $V_{\text{post}} = p\,S_{\text{post}}$ cancels $S_{\text{post}}$ and gives
\[
  p\,S_{\text{exist}} = V_{\text{pre}} - \theta_{\text{pool}} V_{\text{post}},
\]
proving~\eqref{eq:effective-pps} and~\eqref{eq:effective-pre-money}. Because the investor's fraction $\alpha = I / V_{\text{post}}$ is independent of $\theta_{\text{pool}}$, $100\%$ of $\theta_{\text{pool}}$ is deducted from existing holders' post-round share.
\end{proof}

Table~\ref{tab:option-pool-shuffle} and Figure~\ref{fig:option-pool-shuffle} display the exact numbers on $S_{\text{exist}} = 8{,}000{,}000$ founder shares for $V_{\text{pre}} = \text{USD } 16{,}000{,}000$, $I = \text{USD } 4{,}000{,}000$ ($V_{\text{post}} = \text{USD } 20{,}000{,}000$), and $\theta_{\text{pool}} = 10\%$ under two structures: the standard term-sheet \textbf{pre-money pool} ($p = \text{USD } 14{,}000{,}000 / 8{,}000{,}000 = \text{USD } 1.7500$) versus creating the $10\%$ pool \textbf{post-money} (pricing the investment first at $\text{USD } 16{,}000{,}000 / 8{,}000{,}000 = \text{USD } 2.0000$ per share, so founders and investor share the $10\%$ pool dilution proportionally by factor $0.90$).

\begin{table}[htbp]
\centering
\caption{Pre-money option pool shuffle vs.\ post-money pool sharing (\$16M stated pre-money, \$4M investment, $10\%$ pool).}
\label{tab:option-pool-shuffle}
\small
\begin{tabular}{@{}lrrrr@{}}
\toprule
\textbf{Metric / Stakeholder} & \multicolumn{2}{c}{\textbf{Pre-Money Pool (Standard)}} & \multicolumn{2}{c}{\textbf{Post-Money Pool (Shared)}} \\
\cmidrule(lr){2-3} \cmidrule(l){4-5}
 & \textbf{Shares} & \textbf{Ownership} & \textbf{Shares} & \textbf{Ownership} \\
\midrule
Founders (existing shares) & $8{,}000{,}000$ & $70.00\%$ & $8{,}000{,}000$ & $72.00\%$ \\
Unallocated Option Pool    & $1{,}142{,}857$ & $10.00\%$ & $1{,}111{,}111$ & $10.00\%$ \\
Series~A Investor          & $2{,}285{,}714$ & $20.00\%$ & $2{,}000{,}000$ & $18.00\%$ \\
\midrule
\textbf{Total Post-Round ($S_{\text{post}}$)} & $\mathbf{11{,}428{,}571}$ & $\mathbf{100.00\%}$ & $\mathbf{11{,}111{,}111}$ & $\mathbf{100.00\%}$ \\
Investor entry price per share & \multicolumn{2}{r}{\$1.7500} & \multicolumn{2}{r}{\$2.0000 (pre-pool)} \\
Implied value of founder shares & \multicolumn{2}{r}{\$14{,}000{,}000} & \multicolumn{2}{r}{\$14{,}400{,}000} \\
\bottomrule
\end{tabular}
\end{table}

\begin{figure}[htbp]
\centering
\begin{tikzpicture}[x=1cm, y=1cm, font=\small]
  \node[anchor=west, font=\small\bfseries] at (0, 3.35) {Headline Term Sheet: \$16M Pre-Money + \$4M Investment = \$20M Post-Money};

  % Bar 1: Naive reading (No Pool Cut)
  \node[anchor=east, align=right, font=\footnotesize\bfseries] at (2.7, 2.45) {Naive Reading\\[-1pt]\normalfont\scriptsize(No Pool Cut)};
  \fill[sourceblue!80] (2.85, 2.15) rectangle (9.89, 2.75);
  \fill[targetgreen!80] (9.89, 2.15) rectangle (11.65, 2.75);
  \draw[black!60] (2.85, 2.15) rectangle (11.65, 2.75);
  \draw[black!60] (9.89, 2.15) -- (9.89, 2.75);
  \node[white, font=\footnotesize\bfseries] at (6.37, 2.45) {Founders: 80.0\% (\$16.0M, \$2.0000/sh)};
  \node[white, font=\scriptsize\bfseries] at (10.77, 2.45) {Inv: 20\%};

  % Bar 2: Actual Pre-Money Option Pool Shuffle
  \node[anchor=east, align=right, font=\footnotesize\bfseries] at (2.7, 1.35) {Pre-Money Pool\\[-1pt]\normalfont\scriptsize(Effective \$14M Pre)};
  \fill[sourceblue!80] (2.85, 1.05) rectangle (9.01, 1.65);
  \fill[quotientorange!85] (9.01, 1.05) rectangle (9.89, 1.65);
  \fill[targetgreen!80] (9.89, 1.05) rectangle (11.65, 1.65);
  \draw[black!60] (2.85, 1.05) rectangle (11.65, 1.65);
  \draw[black!60] (9.01, 1.05) -- (9.01, 1.65);
  \draw[black!60] (9.89, 1.05) -- (9.89, 1.65);
  \node[white, font=\footnotesize\bfseries] at (5.93, 1.35) {Founders: 70.0\% (\$14.0M, \$1.7500/sh)};
  \node[white, font=\scriptsize\bfseries] at (9.45, 1.35) {10\%};
  \node[white, font=\scriptsize\bfseries] at (10.77, 1.35) {Inv: 20\%};

  % Bar 3: Post-Money Option Pool
  \node[anchor=east, align=right, font=\footnotesize\bfseries] at (2.7, 0.25) {Post-Money Pool\\[-1pt]\normalfont\scriptsize(Shared Dilution)};
  \fill[sourceblue!80] (2.85, -0.05) rectangle (9.186, 0.55);
  \fill[quotientorange!85] (9.186, -0.05) rectangle (10.066, 0.55);
  \fill[targetgreen!80] (10.066, -0.05) rectangle (11.65, 0.55);
  \draw[black!60] (2.85, -0.05) rectangle (11.65, 0.55);
  \draw[black!60] (9.186, -0.05) -- (9.186, 0.55);
  \draw[black!60] (10.066, -0.05) -- (10.066, 0.55);
  \node[white, font=\footnotesize\bfseries] at (6.01, 0.25) {Founders: 72.0\% (\$14.4M)};
  \node[white, font=\scriptsize\bfseries] at (9.626, 0.25) {10\%};
  \node[white, font=\scriptsize\bfseries] at (10.858, 0.25) {Inv: 18\%};

  % Brace highlighting the $16M stated pre-money = $14M founders + $2M pool
  \draw[decorate, decoration={brace, mirror, amplitude=4pt}, thick, trapred]
    (2.85, -0.22) -- (9.89, -0.22)
    node[midway, below=4pt, font=\scriptsize\bfseries, text=trapred]
    {In row 2, ``\$16M Pre-Money'' = \$14M Founder Value + \$2M Pool Value (100\% founder-funded)};
\end{tikzpicture}
\caption{Placing a $10\%$ post-money option pool inside a \$16M pre-money valuation reduces the effective valuation of existing shares to \$14M (\$1.7500/share vs.\ \$2.0000/share).}
\label{fig:option-pool-shuffle}
\end{figure}

\begin{happeningbox}
Trace what happens in Series~B if that $10\%$ Series~A option pool was oversized. Suppose the company only needed $4\%$ of equity to hire every engineer between Series~A and Series~B. The leftover $6\%$ of unallocated options already sits on the cap table when Series~B arrives, so the Series~B investor requires little or no pool expansion. Who paid for that $6\%$ reserve? The founders alone, at Series~A. Who enjoys reduced dilution from it at Series~B? Both the founders \emph{and} the Series~A investor. An oversized pre-money option pool transfers ownership directly from founders to Series~A investors. The constructive counter-move in any term-sheet negotiation is a bottom-up 18-month hiring plan listing each planned hire and target grant, sizing $\theta_{\text{pool}}$ to real hires rather than a round placeholder.
\end{happeningbox}

\section{Convertible Notes and SAFEs}
\label{sec:convertible-notes-and-safes}

At the pre-seed and seed stages -- when two founders have a repository, three design-partner customers, and uncertain pricing -- negotiating a 60-page Amended and Restated Certificate of Incorporation plus Investor Rights, Voting, and Right-of-First-Refusal agreements creates heavy legal friction. Early seed rounds therefore routinely defer formal share pricing by issuing \textbf{convertible instruments}: contracts in which the investor wires cash $I$ today in exchange for the right to receive shares of Preferred Stock automatically when the company later closes a priced equity financing (typically Series~A).

\subsection{Convertible Debt vs.\ the SAFE}

Two instruments dominate early convertible seed rounds:
\begin{itemize}
  \item \textbf{Convertible Promissory Note}: a debt loan carrying principal $I$, an annual interest rate $r$ (typically $4\%$ to $8\%$, accruing into a larger converting balance $I(1 + rt)$ after $t$ years rather than paid in cash), and a \textbf{maturity date} (usually 18 to 24 months after signing). If no priced financing has occurred by maturity, the note legally comes due in cash -- giving noteholders default leverage to renegotiate terms or threaten insolvency if they decline to extend.
  \item \textbf{SAFE (Simple Agreement for Future Equity)}: a five-page equity contract with \textbf{zero interest rate} ($r = 0$) and \textbf{no maturity date}. A SAFE is not debt; it never comes due in cash while the company continues operating, remaining outstanding until a priced equity round, an acquisition, or a dissolution occurs.
\end{itemize}

\subsection{Valuation Caps, Discount Rates, and Shadow Preferred}

Suppose an angel wires \$500{,}000 on a SAFE today, and 18 months later your engineering execution raises the Series~A valuation to \$20{,}000{,}000. If the angel's \$500{,}000 converted at the exact same share price $p_{\text{Series A}}$ paid by the Series~A investor who waited until the product was de-risked, the angel would receive zero compensation for bearing Day~1 risk. Convertibles reward early risk through one or both of two pricing rules:

\begin{enumerate}
  \item \textbf{Discount Rate} $d \in (0, 1)$ (commonly $d = 0.20$, written as conversion multiplier $1 - d = 0.80$): the SAFE converts at $(1 - d)\,p_{\text{Series A}}$ per share.
  \item \textbf{Valuation Cap} $C_{\text{cap}}$: an upper bound on the company valuation used to compute the SAFE's conversion price. Against the SAFE's defined capitalization denominator $S_{\text{cap}}$, the capped conversion price is $C_{\text{cap}} / S_{\text{cap}}$.
\end{enumerate}

When a SAFE carries both a valuation cap and a discount, the investor receives whichever rule produces the \emph{lower} price per share (and therefore \emph{more} shares):
\begin{equation}\label{eq:safe-conversion-price}
  p_{\text{SAFE}} = \min\!\left(\frac{C_{\text{cap}}}{S_{\text{cap}}},\; (1 - d)\,p_{\text{Series A}}\right),
  \qquad
  S_{\text{SAFE}} = \frac{I}{p_{\text{SAFE}}}.
\end{equation}

Notice an essential charter detail called \textbf{Shadow Preferred} (or \textbf{Sub-Series Preferred}, such as Series~A-1 and Series~A-2). Suppose an angel invests \$500{,}000 at a \$5{,}000{,}000 cap and converts during a \$20{,}000{,}000 Series~A round at one-quarter of $p_{\text{Series A}}$, receiving four times as many shares per dollar as the Series~A lead. If the company issued the angel standard Series~A shares carrying a liquidation preference equal to $p_{\text{Series A}}$ per share, that \$500{,}000 check would suddenly carry $4 \times 500{,}000 = 2{,}000{,}000$ dollars (\$2{,}000{,}000) of liquidation preference -- an unintended $4\times$ preference overhang (see Chapter~\ref{ch:founder-contracts}). Standard SAFE conversion prevents this distortion by issuing SAFE holders a separate shadow sub-series (Series~A-1 Preferred) identical to Series~A except that its per-share liquidation preference and initial conversion price equal the angel's actual $p_{\text{SAFE}}$, keeping total liquidation preference equal to actual cash invested (\$500{,}000).

\subsection{Pre-Money SAFEs vs.\ Stacked YC Post-Money SAFEs}

Everything in equation~\eqref{eq:safe-conversion-price} turns on a single contractual definition: which shares sit inside the cap denominator $S_{\text{cap}}$?

Under the original 2013 \textbf{Pre-Money SAFE}, $S_{\text{cap}}$ included existing founder shares and the option pool (including the new Series~A option pool increase), while \emph{excluding} other SAFEs converting in the same financing. Consequently, each additional seed check diluted earlier SAFE holders alongside the founders, and neither side knew exact ownership percentages until the Series~A term sheet arrived.

In 2018, Y Combinator introduced the \textbf{Post-Money SAFE}. Its rule is exact: the SAFE's capitalization denominator $S_{\text{cap}}$ (titled \emph{Company Capitalization} in the contract) \emph{includes} all existing shares, existing options, and \textbf{all converting SAFEs and convertible notes}, while \emph{excluding} the new Series~A investor shares and the new Series~A option pool increase. Therefore, a Post-Money SAFE of cash $I_k$ on post-money valuation cap $C_k$ (when converting via its cap) locks in an exact ownership fraction
\[
  \omega_k = \frac{I_k}{C_k}
\]
of the company's total pre-Series-A share count $S_{\text{pre-A}}$ immediately prior to the Series~A new-money and option-pool issuance.

Let us derive the exact share counts when a startup stacks $m$ Post-Money SAFEs before Series~A. Suppose existing common and option holders start with $S_0$ shares, and the company signs $m$ Post-Money SAFEs with ownership fractions $\omega_1, \dots, \omega_m$ satisfying $\Omega = \sum_{k=1}^m \omega_k < 1$. By contract, SAFE~$k$ receives $S_{\text{SAFE},k}$ shares satisfying
\[
  \frac{S_{\text{SAFE},k}}{S_{\text{pre-A}}} = \omega_k \qquad (k = 1, \dots, m),
  \quad\text{where}\quad
  S_{\text{pre-A}} = S_0 + \sum_{j=1}^m S_{\text{SAFE},j}.
\]
Summing $S_{\text{SAFE},k} = \omega_k S_{\text{pre-A}}$ across $k = 1, \dots, m$ yields $\sum_{k=1}^m S_{\text{SAFE},k} = \Omega\,S_{\text{pre-A}}$, so
\begin{equation}\label{eq:stacked-post-money-safes}
  S_0 = (1 - \Omega)\,S_{\text{pre-A}}
  \implies
  S_{\text{pre-A}} = \frac{S_0}{1 - \sum_{k=1}^m \omega_k},
  \qquad
  S_{\text{SAFE},k} = \frac{\omega_k}{1 - \sum_{j=1}^m \omega_j}\,S_0.
\end{equation}

Inspect equation~\eqref{eq:stacked-post-money-safes}. If you raise \$500{,}000 on a \$5{,}000{,}000 Post-Money SAFE ($\omega_1 = 10.0\%$) and later raise \$1{,}500{,}000 on a \$10{,}000{,}000 Post-Money SAFE ($\omega_2 = 15.0\%$), SAFE~2 does \emph{not} dilute SAFE~1 prior to Series~A. Together they claim a locked additive slice $\Omega = 10.0\% + 15.0\% = 25.0\%$ of $S_{\text{pre-A}}$, compressing the founders to $75.0\%$ before the Series~A investor or the new Series~A option pool takes a single share. Only when Series~A closes do the Series~A investor and new Series~A option pool dilute both the founders and all converted SAFE holders proportionally.

\begin{trapbox}
\textbf{Stacking Low-Cap SAFEs Without Modeling Conversion.}
Because signing a Post-Money SAFE requires no stockholder vote, no charter amendment, and no updated cap table exhibit on page~1, founders sometimes collect seed checks casually over two years: \$400{,}000 from angels on a \$4M post-money cap ($10.0\%$), \$600{,}000 from an accelerator syndicate on a \$5M cap ($12.0\%$), \$1{,}000{,}000 from a micro-VC on an \$8M cap ($12.5\%$), and a \$1{,}500{,}000 seed extension on a \$10M cap ($15.0\%$).

Sum the locked slices: $\Omega = 10.0\% + 12.0\% + 12.5\% + 15.0\% = 49.5\%$ of the pre-Series-A company has already been sold. Once the company raises a standard Series~A selling $20\%$ to the Series~A lead plus a $10\%$ post-money option pool (leaving $70\%$ for all pre-Series-A holders combined), the founders' combined ownership drops to
\[
  (1 - 0.495) \times 70.0\% = 35.35\%
\]
-- giving away nearly $65\%$ of the company by the end of Series~A (and $49.5\%$ before Series~A even opened). Model equation~\eqref{eq:stacked-post-money-safes} in a live cap-table spreadsheet before signing every convertible instrument.
\end{trapbox}

\section{Complete Worked Cap Table Evolution}
\label{sec:worked-cap-table-evolution}

Let us trace one startup from incorporation through two Post-Money SAFEs and a priced Series~A round, recording exact share counts, prices per share to four decimal places, and ownership percentages for every stakeholder.

\subsection{Stage 0: Incorporation}

Two co-founders, Founder~A and Founder~B, incorporate a Delaware C-Corporation with $S_{\text{auth}} = 10{,}000{,}000$ authorized shares of Common Stock at \$0.00001 par value. Each founder buys $4{,}000{,}000$ common shares for \$40.00 cash ($S_0 = 8{,}000{,}000$ shares issued and outstanding, $50.00\%$ each; zero option pool).

\subsection{Stage 1: Two Post-Money Seed SAFEs}

Over its first 14 months, the startup signs two YC Post-Money SAFEs (cap only, zero discount):
\begin{itemize}
  \item \textbf{SAFE~1 (Pre-Seed Angel)}: \$500{,}000 on a \$5{,}000{,}000 Post-Money Valuation Cap, locking in $\omega_1 = 500{,}000 / 5{,}000{,}000 = 10.00\%$ of $S_{\text{pre-A}}$.
  \item \textbf{SAFE~2 (Seed Fund)}: \$1{,}500{,}000 on a \$10{,}000{,}000 Post-Money Valuation Cap, locking in $\omega_2 = 1{,}500{,}000 / 10{,}000{,}000 = 15.00\%$ of $S_{\text{pre-A}}$.
\end{itemize}
Combined SAFE ownership prior to Series~A is $\Omega = 0.10 + 0.15 = 0.25$ ($25.00\%$).

\subsection{Stage 2: Priced Series A With a 10\% Post-Money Option Pool}

At Month~20, a venture fund signs a Series~A term sheet investing $I_A = \text{USD } 5{,}000{,}000$ at a stated Pre-Money Valuation of $V_{\text{pre}} = \text{USD } 20{,}000{,}000$ ($V_{\text{post}} = \text{USD } 25{,}000{,}000$, so $\alpha_A = 5 / 25 = 20.00\%$), including an unallocated Post-Money Option Pool equal to $\theta_{\text{pool}} = 10.00\%$ of post-Series-A fully diluted shares (\$2{,}500{,}000 value) inside the pre-money capitalization alongside the converting SAFEs.

Execute the three-step closing calculation:

\textbf{Step 1 (Convert the Post-Money SAFEs).} By equation~\eqref{eq:stacked-post-money-safes} with $S_0 = 8{,}000{,}000$ and $1 - \Omega = 0.75$:
\[
  S_{\text{pre-A}} = \frac{8{,}000{,}000}{0.75} = 10{,}666{,}667 \text{ shares}
\]
(rounding whole shares to the nearest integer):
\begin{align*}
  S_{\text{SAFE,1}} &= 0.10 \times \frac{8{,}000{,}000}{0.75} = 1{,}066{,}667 \text{ shares of Series~A-1 Preferred}, \\
  S_{\text{SAFE,2}} &= 0.15 \times \frac{8{,}000{,}000}{0.75} = 1{,}600{,}000 \text{ shares of Series~A-2 Preferred}.
\end{align*}
The effective conversion prices per share paid by the SAFE holders are
\begin{align*}
  p_{\text{SAFE,1}} &= \frac{\text{USD } 500{,}000}{1{,}066{,}667} = \text{USD } 0.4688 \quad \Bigl(\text{exact } \frac{15}{32} = 0.46875\Bigr), \\
  p_{\text{SAFE,2}} &= \frac{\text{USD } 1{,}500{,}000}{1{,}600{,}000} = \text{USD } 0.9375 \quad \Bigl(\text{exact } \frac{15}{16}\Bigr).
\end{align*}

\textbf{Step 2 (Compute the Series~A Share Price).} Applying Proposition~\ref{prop:option-pool-shuffle} to the $S_{\text{pre-A}} = 10{,}666{,}667$ pre-Series-A shares, the Series~A investor ($20.00\%$) and new unallocated option pool ($10.00\%$) leave $70.00\%$ of post-Series-A equity -- worth $V_{\text{eff}} = \text{USD } 20{,}000{,}000 - \text{USD } 2{,}500{,}000 = \text{USD } 17{,}500{,}000$ -- for the combined Founders + SAFE holders:
\[
  p_A = \frac{\text{USD } 17{,}500{,}000}{10{,}666{,}667} = \text{USD } 1.6406 \text{ per share}
  \quad \Bigl(\text{exact } \frac{105}{64} = 1.640625\Bigr).
\]

\textbf{Step 3 (Issue Series~A Preferred and Option Pool Shares, and Amend Charter).} Dividing dollar allocations by exact $p_A = 105/64 = \text{USD } 1.640625$ yields integer share counts:
\begin{align*}
  S_{\text{Series A}} &= \frac{\text{USD } 5{,}000{,}000}{\text{USD } 1.640625} = 3{,}047{,}619 \text{ shares of Series~A Preferred}, \\
  S_{\text{pool}} &= \frac{\text{USD } 2{,}500{,}000}{\text{USD } 1.640625} = 1{,}523{,}810 \text{ shares reserved in Option Pool}.
\end{align*}
Total post-Series-A fully diluted shares reach $S_{\text{post}} = 10{,}666{,}667 + 1{,}523{,}810 + 3{,}047{,}619 = 15{,}238{,}096$, exceeding the original $10{,}000{,}000$ Incorporation charter limit. At closing, the Amended and Restated Certificate of Incorporation expands authorized Common Stock (say to $20{,}000{,}000$ shares, covering common, options, and full preferred conversion) and authorizes $5{,}714{,}286$ shares of Preferred Stock across Series~A, Series~A-1, and Series~A-2. Table~\ref{tab:cap-table-evolution} and Figure~\ref{fig:cap-table-evolution} summarize all three stages.

\begin{table}[htbp]
\centering
\caption{Complete cap table evolution from Incorporation through two Post-Money SAFEs and Series~A (\$5M on \$20M Pre-Money, $10\%$ Post-Money Pool).}
\label{tab:cap-table-evolution}
\footnotesize
\setlength{\tabcolsep}{3.4pt}
\begin{tabular}{@{}llrrrrrr@{}}
\toprule
 & & \multicolumn{2}{c}{\textbf{Incorporation}} & \multicolumn{2}{c}{\textbf{Pre-A (SAFEs)}} & \multicolumn{2}{c}{\textbf{Post-Series A}} \\
\cmidrule(lr){3-4} \cmidrule(lr){5-6} \cmidrule(l){7-8}
\textbf{Stakeholder} & \textbf{Security Class} & \textbf{Shares} & \textbf{\% FD} & \textbf{Shares} & \textbf{\% FD} & \textbf{Shares} & \textbf{\% FD} \\
\midrule
Founder A & Common & $4{,}000{,}000$ & $50.00\%$ & $4{,}000{,}000$ & $37.50\%$ & $4{,}000{,}000$ & $26.25\%$ \\
Founder B & Common & $4{,}000{,}000$ & $50.00\%$ & $4{,}000{,}000$ & $37.50\%$ & $4{,}000{,}000$ & $26.25\%$ \\
SAFE 1 (\$500k @ \$5M) & Series A-1 Pref & $0$ & $0.00\%$ & $1{,}066{,}667$ & $10.00\%$ & $1{,}066{,}667$ & $7.00\%$ \\
SAFE 2 (\$1.5M @ \$10M) & Series A-2 Pref & $0$ & $0.00\%$ & $1{,}600{,}000$ & $15.00\%$ & $1{,}600{,}000$ & $10.50\%$ \\
Employee Option Pool & Common Options & $0$ & $0.00\%$ & $0$ & $0.00\%$ & $1{,}523{,}810$ & $10.00\%$ \\
Series A Lead (\$5M) & Series A Pref & $0$ & $0.00\%$ & $0$ & $0.00\%$ & $3{,}047{,}619$ & $20.00\%$ \\
\midrule
\textbf{Total Fully Diluted} & & $\mathbf{8{,}000{,}000}$ & $\mathbf{100.00\%}$ & $\mathbf{10{,}666{,}667}$ & $\mathbf{100.00\%}$ & $\mathbf{15{,}238{,}096}$ & $\mathbf{100.00\%}$ \\
Entry / Round Price & & \multicolumn{2}{r}{\$0.00001/sh} & \multicolumn{2}{r}{\$0.4688 / \$0.9375} & \multicolumn{2}{r}{\$1.6406/sh} \\
Post-Round Valuation & & \multicolumn{2}{r}{\$80 (par)} & \multicolumn{2}{r}{--} & \multicolumn{2}{r}{\$25{,}000{,}000} \\
\bottomrule
\end{tabular}
\end{table}

\begin{figure}[htbp]
\centering
\begin{tikzpicture}[x=1cm, y=1cm, font=\small]
  \node[anchor=west, font=\small\bfseries] at (0, 3.65) {Ownership Dilution Across Stages (100\% Fully Diluted Width)};

  % Stage 0: Incorporation (50% + 50%)
  \node[anchor=east, align=right, font=\footnotesize\bfseries] at (2.45, 2.75) {Stage 0\\[-1pt]\normalfont\scriptsize Incorporation\\[-2pt]\normalfont\scriptsize (8.00M sh)};
  \fill[sourceblue!85] (2.60, 2.40) rectangle (7.35, 3.10);
  \fill[sourceblue!60] (7.35, 2.40) rectangle (12.10, 3.10);
  \draw[black!60] (2.60, 2.40) rectangle (12.10, 3.10);
  \draw[white, thick] (7.35, 2.40) -- (7.35, 3.10);
  \node[white, font=\footnotesize\bfseries] at (4.975, 2.75) {Founder A: 50.00\%};
  \node[white, font=\footnotesize\bfseries] at (9.725, 2.75) {Founder B: 50.00\%};

  % Stage 1: SAFE Conversion Pre-A (37.5% + 37.5% + 10.0% + 15.0%)
  \node[anchor=east, align=right, font=\footnotesize\bfseries] at (2.45, 1.55) {Stage 1\\[-1pt]\normalfont\scriptsize SAFEs Convert\\[-2pt]\normalfont\scriptsize (10.67M sh)};
  \fill[sourceblue!85] (2.60, 1.20) rectangle (6.1625, 1.90);
  \fill[sourceblue!60] (6.1625, 1.20) rectangle (9.725, 1.90);
  \fill[structurepurple!80] (9.725, 1.20) rectangle (10.675, 1.90);
  \fill[structurepurple!55] (10.675, 1.20) rectangle (12.10, 1.90);
  \draw[black!60] (2.60, 1.20) rectangle (12.10, 1.90);
  \draw[white, thick] (6.1625, 1.20) -- (6.1625, 1.90);
  \draw[white, thick] (9.725, 1.20) -- (9.725, 1.90);
  \draw[white, thick] (10.675, 1.20) -- (10.675, 1.90);
  \node[white, font=\scriptsize\bfseries] at (4.38, 1.55) {Founder A: 37.5\%};
  \node[white, font=\scriptsize\bfseries] at (7.94, 1.55) {Founder B: 37.5\%};
  \node[white, font=\scriptsize\bfseries] at (10.20, 1.55) {S1: 10\%};
  \node[white, font=\scriptsize\bfseries] at (11.39, 1.55) {S2: 15\%};

  % Connector lines showing proportional 0.70 compression from Stage 1 to Stage 2
  \draw[dashed, black!45] (9.725, 1.20) -- (7.5875, 0.70);
  \draw[dashed, black!45] (12.10, 1.20) -- (9.25, 0.70);

  % Stage 2: Post-Series A (26.25% + 26.25% + 7.0% + 10.5% + 10.0% + 20.0%)
  \node[anchor=east, align=right, font=\footnotesize\bfseries] at (2.45, 0.35) {Stage 2\\[-1pt]\normalfont\scriptsize Post-Series A\\[-2pt]\normalfont\scriptsize (15.24M sh)};
  \fill[sourceblue!85] (2.60, 0.00) rectangle (5.09375, 0.70);
  \fill[sourceblue!60] (5.09375, 0.00) rectangle (7.5875, 0.70);
  \fill[structurepurple!80] (7.5875, 0.00) rectangle (8.2525, 0.70);
  \fill[structurepurple!55] (8.2525, 0.00) rectangle (9.25, 0.70);
  \fill[quotientorange!85] (9.25, 0.00) rectangle (10.20, 0.70);
  \fill[targetgreen!80] (10.20, 0.00) rectangle (12.10, 0.70);
  \draw[black!60] (2.60, 0.00) rectangle (12.10, 0.70);
  \foreach \x in {5.09375, 7.5875, 8.2525, 9.25, 10.20} {
    \draw[white, thick] (\x, 0.00) -- (\x, 0.70);
  }
  \node[white, font=\scriptsize\bfseries] at (3.85, 0.35) {A: 26.25\%};
  \node[white, font=\scriptsize\bfseries] at (6.34, 0.35) {B: 26.25\%};
  \node[white, font=\tiny\bfseries] at (7.92, 0.35) {7\%};
  \node[white, font=\scriptsize\bfseries] at (8.75, 0.35) {10.5\%};
  \node[white, font=\scriptsize\bfseries] at (9.725, 0.35) {Pool 10\%};
  \node[white, font=\scriptsize\bfseries] at (11.15, 0.35) {Ser A: 20\%};
\end{tikzpicture}
\caption{Two-phase dilution across Table~\ref{tab:cap-table-evolution}: Post-Money SAFEs~1 and~2 first absorb $25.00\%$ purely from the founders; Series~A ($20.00\%$) and the new Option Pool ($10.00\%$) then compress both founders and converted SAFEs by a uniform factor of $0.70$.}
\label{fig:cap-table-evolution}
\end{figure}

Check the multiplicative structure in Table~\ref{tab:cap-table-evolution}: each founder's $50.00\%$ initial stake is multiplied first by $(1 - 0.25) = 0.75$ from the two Post-Money SAFEs ($37.50\%$), and then by $(1 - 0.20 - 0.10) = 0.70$ from Series~A and the new option pool ($0.375 \times 0.70 = 26.25\%$). Meanwhile, SAFE~1's $10.00\%$ and SAFE~2's $15.00\%$ also compress by $0.70$ to $7.00\%$ and $10.50\%$. At $p_A = \text{USD } 1.640625$ per share, each founder's $4{,}000{,}000$ shares carry an implied paper value of \$6{,}562{,}500, SAFE~1's \$500{,}000 check is marked at \$1{,}750{,}000 ($3.50\times$), and SAFE~2's \$1{,}500{,}000 check is marked at \$2{,}625{,}000 ($1.75\times$). Whether those paper marks turn into actual cash upon a sale depends on the preferred liquidation waterfall, vesting rules, and founder compensation contracts we develop in Chapter~\ref{ch:founder-contracts}.

\section*{Exercises}
\addcontentsline{toc}{section}{Exercises}

\begin{exercise}[Default Alive vs.\ Default Dead Cash Calculus]
A developer-tools startup holds $C_0 = \text{USD } 1{,}200{,}000$ in cash, incurs fixed monthly operating expenses $E_0 = \text{USD } 160{,}000$, and earns initial monthly revenue $R_0 = \text{USD } 75{,}000$ at $m = 80\%$ gross margin. Compute: (a) initial monthly net burn $B_0$ and static runway $\tau_{\text{static}}$; (b) break-even month $t^*$ and cumulative cash required $\Delta C_{\text{req}}$ under equation~\eqref{eq:default-alive-cash} if monthly continuous organic growth is $g = 0.05$ ($5\%$ per month), stating whether the company is Default Alive and its minimum cash balance $C_{\min}$; and (c) the minimum continuous monthly growth rate $g_{\min}$ required for the startup to be Default Alive without cutting $E_0$.
\end{exercise}

\begin{exercise}[Authorized vs.\ Outstanding vs.\ Fully Diluted]
A corporation's charter authorizes $15{,}000{,}000$ shares of Common Stock. Three co-founders hold $2{,}500{,}000$ shares each. The board has reserved $1{,}500{,}000$ shares in an Equity Incentive Plan, of which $600{,}000$ options have been granted to early engineers and $900{,}000$ remain unallocated. Compute: (a) Issued and Outstanding shares $S_{\text{out}}$; (b) Fully Diluted shares $S_{\text{FD}}$; (c) unissued and unreserved authorized shares; and (d) one founder's ownership percentage on both a voting outstanding basis and a fully diluted basis.
\end{exercise}

\begin{exercise}[Pricing a Clean Seed Round vs.\ an Option Pool Shuffle]
A startup has $9{,}000{,}000$ issued founder shares and $0$ reserved options. Investor~X offers \$3{,}000{,}000 at a \$12{,}000{,}000 Pre-Money Valuation ($V_{\text{post}} = \text{USD } 15{,}000{,}000$) with a $12\%$ Post-Money Option Pool created inside the pre-money capitalization. Investor~Y offers \$3{,}000{,}000 at an \$11{,}000{,}000 Pre-Money Valuation ($V_{\text{post}} = \text{USD } 14{,}000{,}000$) requiring only a $5\%$ Post-Money Option Pool inside the pre-money capitalization. For both offers, compute the effective pre-money valuation $V_{\text{eff}}$ of the founders' shares, the round price per share to four decimal places, and the founders' post-round ownership percentage. Which offer values the founders' existing shares higher?
\end{exercise}

\begin{exercise}[Cap vs.\ Discount Crossover]
An angel invests \$400{,}000 in a Post-Money SAFE carrying both an \$8{,}000{,}000 Post-Money Valuation Cap and a $20\%$ Discount ($1 - d = 0.80$). Assume no other SAFEs exist and existing founders hold $8{,}000{,}000$ shares, with no new Series~A option pool required. At what Series~A pre-money valuation $V_{\text{pre}}$ do the valuation cap and the $20\%$ discount produce the exact same conversion price per share? Compute the angel's converted share count and conversion price per share (to four decimal places) if Series~A raises \$4{,}000{,}000 at (a) $V_{\text{pre}} = \text{USD } 8{,}500{,}000$ and (b) $V_{\text{pre}} = \text{USD } 16{,}000{,}000$.
\end{exercise}

\begin{exercise}[Unmodeled Stacked SAFEs]
Three founders start with $3{,}000{,}000$ shares each ($S_0 = 9{,}000{,}000$). Over two years they sign three Post-Money SAFEs: \$300{,}000 on a \$3{,}000{,}000 post-money cap, \$800{,}000 on an \$8{,}000{,}000 post-money cap, and \$2{,}000{,}000 on a \$12{,}500{,}000 post-money cap. They then raise a Series~A of \$6{,}000{,}000 at a \$24{,}000{,}000 Pre-Money Valuation ($V_{\text{post}} = \text{USD } 30{,}000{,}000$) including an $8\%$ Post-Money Option Pool inside the pre-money capitalization. Compute $S_{\text{pre-A}}$, the Series~A price per share $p_A$ to four decimal places, the post-Series-A share count for every stakeholder, and each individual founder's post-Series-A ownership percentage.
\end{exercise}
